% 原书第 36—37 页；3.3。
\section{反事实}
\label{sec:ch03-counterfactuals}

对结构因果模型还可以作另一种修改：改变它的所有噪声分布。这种变化可以由观测引起，并使我们能够回答\textbf{反事实（counterfactual）}问题。为说明这一点，请考虑下面这个假想情境。

\begin{example}[眼病]
某种眼病有一种相当有效的治疗方法。对所有患者中的 $99\%$ 而言，治疗能够发挥作用，使其痊愈（$B=0$）；如果不治疗，这些患者就会在一天内失明（$B=1$）。对余下的 $1\%$ 而言，治疗的作用恰好相反，会使其在一天内失明（$B=1$）；若不治疗，他们反而会恢复正常视力（$B=0$）。

患者属于哪一类，取决于其是否具有某种罕见状况（$N_B=1$）。医生并不知道患者是否具有这种状况，因此，其是否给予治疗（$T=1$）的决定与 $N_B$ 相互独立。我们用噪声变量 $N_T$ 表示这一决定。

假设其背后的结构因果模型为
\begin{equation}
 \mathfrak{C}:\quad
 \begin{aligned}
 T&:=N_T\\
 B&:=T\cdot N_B+(1-T)\cdot(1-N_B),
 \end{aligned}
 \tag{3.5}
\end{equation}
其中 $N_B\sim\operatorname{Ber}(0.01)$ 服从 Bernoulli 分布。注意，相应的因果图为 $T\to B$。

现在，设想一位视力不佳的患者来到医院，在医生给予治疗（$T=1$）后失明了（$B=1$）。我们便可以提出反事实问题：“如果医生当时选择 $T=0$，会发生什么？”令人意外的是，这个问题能够得到回答。根据式 (3.5)，观测到 $B=T=1$ 意味着，对这位患者有 $N_B=1$。有了这一信息，我们就能计算 $\dox{T:=0}$ 的效应。

为此，我们先以观测结果为条件，更新噪声变量的分布。前面已经看到，给定 $B=T=1$，$N_B$ 和 $N_T$ 的分布都会退化为位于 $1$ 的点质量，即 $\delta_1$。这样便得到修改后的结构因果模型：
\begin{equation}
 \mathfrak{C}\mid B=1,T=1:\quad
 \begin{aligned}
 T&:=1\\
 B&:=T\cdot1+(1-T)\cdot(1-1)=T.
 \end{aligned}
 \tag{3.6}
\end{equation}
注意，我们只更新噪声分布；条件化并不改变赋值本身的结构。这里的想法是，物理机制保持不变——在本例中，即导致痊愈或失明的机制保持不变——但对于\emph{这位特定患者}，我们获得了有关先前未知噪声变量的信息。

接下来，计算对这位患者实施 $\dox{T=0}$ 的效应：
\begin{equation}
 \mathfrak{C}\mid B=1,T=1;\dox{T:=0}:\quad
 \begin{aligned}
 T&:=0\\
 B&:=T.
 \end{aligned}
 \tag{3.7}
\end{equation}
显然，该模型导出的分布将全部概率质量都放在 $(0,0)$ 上，因此
\[
 P^{\mathfrak{C}\mid B=1,T=1;\dox{T:=0}}(B=0)=1.
\]
这意味着，如果医生当时没有给予治疗，也就是实施 $\dox{T:=0}$，这位患者就会痊愈（$B=0$）。不过，由于
\[
 \begin{aligned}
 P^{\mathfrak{C};\dox{T:=1}}(B=0)&=0.99,\\
 P^{\mathfrak{C};\dox{T:=0}}(B=0)&=0.01,
 \end{aligned}
\]
我们仍然可以认为，依据当时掌握的知识，医生作出了最优选择。
\end{example}

有意思的是，例 3.4 表明，我们可以利用反事实陈述来否证其背后的因果模型（见第 6.8 节）。设想可以对罕见状况 $N_B$ 进行检测，但需要超过一天才能拿到检测结果。这样就可能出现一种情况：我们得到的反事实陈述与 $N_B$ 的测量结果相矛盾。Pearl［2009，第 220 页，第 (2) 点］也给出了同样的论证。鉴于反事实的科学内容一直存在广泛争论，这里需要强调：上述反事实陈述具有可否证性，是因为噪声变量并非在原则上无法观测，而只是在医生必须作出决定时尚不可观测。
